\section{相关工作}
\label{sec:related_work}

本文的研究位于四条文献脉络的交叉处：LLM 驱动的多智能体系统、面向智能体协作的 benchmark、合作博弈式贡献归因，以及多智能体强化学习中的 credit assignment。与已有工作相比，本文的重点不是提出新的 MAS 架构，也不是仅报告完整系统的成功率，而是系统研究单个智能体的贡献得分如何由拓扑位置、角色、权限、通信行为和任务难度共同决定。

\subsection{LLM 智能体与多智能体系统}

LLM-based autonomous agents 通常将语言模型与记忆、规划、工具调用和环境反馈结合，使模型能够在多轮交互中完成开放式任务。已有综述从 agent 构成、规划机制、工具使用、记忆模块、交互环境和评估方法等角度总结了该方向的发展\cite{wang2024survey_agents,guo2024llm_ma_survey,chen2024llm_mas_survey}。在此基础上，LLM-MAS 进一步把多个具有不同角色、工具权限或通信协议的智能体组织为协作系统，以期通过分工、互评、辩论、路由或标准化流程提升复杂任务求解能力。

早期的多智能体框架多强调角色设定与语言通信。CAMEL 通过 role-playing communicative agents 探索语言智能体之间的自主合作\cite{li2023camel}；AgentVerse 将多个专家智能体组织为动态协作群体，并观察协作过程中涌现的社会行为\cite{chen2023agentverse}；AutoGen 将多智能体会话抽象为可编程框架，使开发者能够定义 agent 之间的对话模式、人类介入和工具调用方式\cite{wu2023autogen}。这些框架说明，MAS 的性能不仅取决于单个模型能力，也取决于 agent 之间如何通信、分工和控制任务流。

软件工程是 LLM-MAS 最具代表性的应用场景之一。MetaGPT 将标准化操作流程编码到多智能体协作中，用 product manager、architect、engineer 等角色模拟软件公司的工作流\cite{hong2023metagpt}；ChatDev 使用 chat chain 组织需求分析、设计、编码和测试阶段\cite{qian2023chatdev}；MapCoder 将程序求解过程拆分为 retrieval/recall、planning、coding 和 debugging 等阶段\cite{islam2024mapcoder}。这些系统为本文的 Architecture Zoo 提供了自然原型：它们包含清晰的角色标签、通信路径、上下游依赖和最终输出权限，适合研究同一模型能力下不同节点的边际贡献。

然而，上述系统通常以完整 MAS 的任务成功率、代码通过率或人工评分作为主要结果。即使观察到多智能体系统优于单智能体，也难以进一步回答：性能提升到底来自 planner、coder、verifier、critic、router，还是来自某个节点拥有最终输出权限或工具权限？本文正是针对这一缺口，试图从系统级评估转向节点级贡献解释。

\subsection{智能体 Benchmark 与多智能体协作评估}

LLM agent benchmark 最初主要关注单个智能体在交互环境中的任务完成能力。AgentBench 提供多个交互式环境，用于评估 LLM 作为 agent 的推理、决策和行动能力\cite{liu2023agentbench}；AgentBoard 强调多轮 agent 评估中的过程分析，指出只看最终成功率会掩盖模型在中间步骤中的进展\cite{ma2024agentboard}；MINT 关注模型在多轮交互中使用工具和利用自然语言反馈的能力\cite{wang2023mint}。这些 benchmark 共同推动了 agent 评估从静态问答走向多轮、工具增强和过程可解释的方向，但它们大多仍以单 agent 或 user--agent--tool 交互为主。

面向 LLM-MAS 的 benchmark 开始显式评估协作、竞争和通信。MultiAgentBench/MARBLE 覆盖多种协作与竞争场景，并将 evaluation 从任务成功扩展到 planning、communication 和 coordination 等指标\cite{zhu2025multiagentbench}。LLM-Coordination 从 pure coordination games 出发，评估 LLM 在环境理解、Theory of Mind 和联合规划中的协调能力\cite{agashe2025llmcoordination}。MindAgent/CuisineWorld 在游戏环境中评估多智能体规划和协作效率\cite{gong2023mindagent}，SOTOPIA 则关注开放式社会互动中的 social intelligence\cite{zhou2024sotopia}。这些 benchmark 说明，多智能体能力并不能由单个最终分数完全刻画，通信过程、协作模式和角色互动本身也是重要评估对象。

近期工作进一步强调角色隔离、私有信息、隐私和协作失败模式。TeamBench 通过操作系统级角色隔离，将 specification access、workspace editing 和 final certification 分配给不同角色，从而暴露 prompt-only role assignment 可能掩盖的越权和虚假协作问题\cite{kim2026teambench}。CoLLAB 将分布式约束优化问题扩展为自然语言多智能体协调环境，使 agent 必须在局部信息和通信约束下优化全局目标\cite{mahmud2025collab}。CalBench 以日程协调为场景，研究 private calendar 下的协调质量、通信效率、公平性和隐私泄露\cite{zou2026calbench}。AgentCollabBench 则从 instruction decay、tracer durability、consensus pollution 和 cross-task leakage 等角度诊断“单个能力强的 agent 组合后为何协作失败”\cite{mazumder2026agentcollabbench}。

这些 benchmark 为本文提供了任务环境和过程指标，但它们通常不把单个 agent 的边际贡献作为核心对象。本文与其互补：我们将这些 benchmark 中的角色、权限、通信和过程指标转化为解释变量，并进一步用 Shapley、LOO、Banzhaf、Myerson 和 Owen 等方法计算 agent-level contribution，从而分析哪些结构性因素使一个 agent 重要。

\subsection{合作博弈、解释性归因与 Agent 贡献}

合作博弈为“个体对集体结果的贡献”提供了经典形式化工具。Shapley value 基于效率、对称性、虚拟玩家和可加性等公理，为每个参与者分配其平均边际贡献\cite{shapley1953value}。Banzhaf index 更强调一个参与者在不同 coalition 中是否成为 pivotal player\cite{banzhaf1965weighted}。Myerson value 将通信图或合作图纳入博弈，使只有图连通或图约束下可行的联盟才能贡献价值\cite{myerson1977graphs}。Owen value 则处理存在先验联盟或层级分组的合作博弈\cite{owen1977values}。在机器学习解释性中，SHAP 将 Shapley 思想用于特征归因\cite{lundberg2017unified}，Data Shapley 将其用于训练样本价值评估\cite{ghorbani2019data}。这些工作证明，合作博弈能够为复杂系统中的组成单元提供原则化贡献分配。

LLM agent 研究近年开始将合作博弈用于模块或智能体贡献分析。CapaBench 将 planning、reasoning、action 和 reflection 等能力模块视为 coalition 成员，用 Shapley value 衡量模块对单个 LLM agent 的边际贡献\cite{yang2025capabench}。An Attribution Framework for Multi-Agent LLMs 将 MAS workflow 形式化为合作博弈，比较 Shapley、Banzhaf 和 leave-one-out 等归因方法，用于识别高贡献或冗余 agent\cite{lu2026attribution}。Agents that Matter 进一步把 agent attribution 参数化为 coalition distribution、removal protocol 和 target metric，并讨论 agent ablation 与 model replacement 所诱导的不同 attribution game\cite{lu2026agents}。HiveMind 在股票交易 MAS 中使用 Shapley value 和 DAG-Shapley 进行贡献引导的在线 prompt 优化，以降低归因计算成本并优化低贡献 agent\cite{xia2026hivemind}。

上述工作与本文最为接近，但研究重点不同。已有 attribution 工作主要回答“哪个 agent 重要”或“如何用贡献得分优化系统”；本文进一步追问“为什么它重要”。具体而言，本文把 Shapley/LOO/Banzhaf/Myerson/Owen 得分作为因变量，将 topology、role、permission、trace behavior 和 task difficulty 作为解释变量，并通过混合效应回归、role-preserving topology rewiring、topology-preserving role swap 和 permission-only intervention 等方法检验结构性假设。因此，本文不仅提供贡献排序，也试图建立 agent importance 的可解释机制。

\subsection{拓扑结构、通信机制与角色/权限因素}

多智能体系统本质上是带有通信、依赖和控制关系的网络。经典网络分析提出了 degree、betweenness、closeness、eigenvector centrality 和 PageRank 等中心性指标，用于刻画节点在信息传播和结构控制中的位置\cite{freeman1977set,bonacich1987power,page1999pagerank,newman2010networks}。在 MAS 中，这些指标对应 supervisor、router、bridge、finalizer 或高扇入/高扇出节点的结构位置。Myerson value 进一步说明，当合作受到图结构约束时，贡献分配本身应显式依赖通信图\cite{myerson1977graphs}。

不过，LLM-MAS 中的节点重要性不一定由拓扑中心性单独决定。一个低 degree 的 finalizer 可能因为拥有最终输出权限而贡献很高；一个高 degree 的 supervisor 可能只是转发信息而没有实质增益；一个 verifier 可能在简单任务上引入延迟或错误批准，却在困难任务上显著减少失败。TeamBench 的角色隔离设置表明，prompt 中声明的角色并不等价于真实权限和真实行为\cite{kim2026teambench}。因此，本文将拓扑因素与角色因素、权限因素和运行时行为因素分开建模，并通过反事实干预区分“结构位置重要”“角色本身重要”和“权限配置重要”三种机制。

通信机制也是 MAS 研究中的重要变量。Communication-centric survey 将 LLM-MAS 的通信划分为系统级通信和系统内部通信，强调通信协议、对象、内容和策略都会影响集体智能\cite{yan2025communication_survey}。MultiAgentBench、CoLLAB、CalBench 和 AgentCollabBench 等 benchmark 也分别从 communication score、局部信息协作、隐私边界和错误传播等角度提供了可观测指标\cite{zhu2025multiagentbench,mahmud2025collab,zou2026calbench,mazumder2026agentcollabbench}。本文吸收这些视角，将 messages sent/received、tokens sent/received、被引用次数、tool-call 成功率、constraint survival 和 role violation 等 trace features 纳入贡献解释模型。

\subsection{多智能体强化学习中的 Credit Assignment}

Credit assignment 是传统多智能体强化学习中的核心问题。在完全合作 MARL 中，系统通常只能观察到全局回报，而学习算法需要估计每个 agent 的动作对团队回报的影响。Difference rewards 通过比较实际团队回报与移除某 agent 影响后的反事实回报，缓解全局奖励难以分配的问题\cite{wolpert2001optimal}。Value Decomposition Networks 将团队价值分解为各个 agent 的局部价值之和\cite{sunehag2018value}；QMIX 使用单调混合网络提高值函数分解的表达能力\cite{rashid2018qmix}；COMA 使用 centralized critic 构造 counterfactual baseline 来估计单个 agent 动作的优势\cite{foerster2018coma}。

这些 MARL 方法与本文共享“从全局结果分解个体贡献”的动机，但问题设置存在明显差异。第一，MARL 通常在训练过程中学习策略和价值函数，而本文面对的是黑箱或半黑箱 LLM agent workflow，重点是部署后诊断和离线归因。第二，MARL 的贡献对象多为动作或策略，本文的贡献对象是带有角色、权限、工具和通信位置的 agent 节点。第三，LLM-MAS 的输出是自然语言、代码、工具调用和文件修改的混合 trace，贡献受到 prompt、上下文、权限和拓扑的共同影响。因此，本文采用合作博弈式 coalition evaluation 和 removal/replacement protocol，而不是直接学习可微分的价值分解模型。

\subsection{本文与已有工作的关系}

表\ref{tab:related_work_positioning} 总结了本文与主要相关方向的区别。总体而言，现有 LLM-MAS 架构工作证明了多角色协作的可行性，benchmark 工作提供了协作任务和过程指标，归因工作开始量化单个 agent 的边际贡献，MARL credit assignment 提供了长期的概念基础。本文的独特定位在于：以 agent-level contribution 为核心因变量，系统检验 topology、role、permission、trace behavior 和 task difficulty 对贡献的解释力，并通过受控反事实干预验证这些因素是否具有稳健影响。

\begin{table}[t]
\centering
\small
\caption{相关工作与本文的定位比较。}
\label{tab:related_work_positioning}
\begin{tabularx}{\textwidth}{L{0.19\textwidth}Y L{0.19\textwidth}L{0.23\textwidth}}
\toprule
研究方向 & 代表工作 & 主要关注 & 与本文的关系 \\
\midrule
LLM-MAS 架构 & CAMEL、AutoGen、AgentVerse、MetaGPT、ChatDev、MapCoder & 构造多角色、多工具或多阶段协作系统，提高复杂任务表现。 & 本文将这些系统抽象为 Architecture Zoo，用于分析节点贡献来源。 \\
多智能体 benchmark & MultiAgentBench、TeamBench、CoLLAB、CalBench、MindAgent、SOTOPIA、AgentCollabBench & 评估协作、竞争、角色隔离、隐私、通信和过程失败。 & 本文复用其任务和过程指标，但进一步进行 agent-level attribution。 \\
合作博弈归因 & Shapley、Banzhaf、Myerson、Owen、SHAP、Data Shapley & 为复杂系统中的组成单元分配边际贡献。 & 本文将这些方法用于 LLM-MAS 节点，并比较不同 attribution protocol。 \\
LLM agent 贡献分析 & CapaBench、Agent That Matters、Agents that Matter、HiveMind & 计算模块或 agent 的贡献，识别冗余或低贡献组件。 & 本文从“贡献计算”进一步推进到“贡献机制解释”。 \\
MARL credit assignment & Difference rewards、VDN、QMIX、COMA & 在全局回报下学习个体动作或策略贡献。 & 本文借鉴反事实思想，但面向黑箱 LLM workflow 的离线诊断。 \\
\bottomrule
\end{tabularx}
\end{table}
